\subsection{Finite-Pool Properties and a Deterministic Coverage Decomposition}
\label{sec:theory-scope}

This section states exactly what follows from the implemented margin sets and
what does not.  The analysis is conditional on the frozen scorer and on all
information available before timestamp \(t\).  It requires no exchangeability
assumption, but consequently its drift and estimation terms are not guaranteed
to be small.

Let \(A_t=a(X_t,Y_t)\) be the margin nonconformity score of a randomly indexed
observed query--label pair from the current timestamp batch.  Conditional on
the revealed past and the unlabeled current batch, define its CDF as
\[
  F_t(u)=\Pr\{A_t\leq u\mid\mathcal{I}_{t^-}\}.
\]
The conditioning information \(\mathcal{I}_{t^-}\) includes the frozen scorer,
past revealed labels, current query identities, and any drift features computed
without current labels.  Let \(a_1,\ldots,a_n\) be the scores in the current
calibration pool, with nonnegative normalized weights
\(\bar w_i\) satisfying \(\sum_i\bar w_i=1\).  If \(F_i\) denotes the
conditional score CDF associated with the source position of \(a_i\), define
the population weighted-history mixture and its empirical counterpart by
\[
  H_{t,w}(u)=\sum_{i=1}^{n}\bar w_iF_i(u),
  \qquad
  \widehat H_{t,w}(u)=\sum_{i=1}^{n}\bar w_i
  \mathbb{1}\{a_i\leq u\}.
\]
For \(p=1-\alpha\), the implemented inclusive weighted threshold is
\[
  q_t=\inf\{u:\widehat H_{t,w}(u)\geq p\}.
\]
Two quantities separate temporal mismatch from finite-pool estimation:
\[
\begin{aligned}
  \delta_t&=\sup_u|F_t(u)-H_{t,w}(u)|,\\
  \epsilon_t&=\sup_u|\widehat H_{t,w}(u)-H_{t,w}(u)|.
\end{aligned}
\]
Because KG score margins can tie, the relevant discretization term is not the
largest individual weight.  It is the largest empirical atom mass
\[
  \eta_t=\max_v\sum_{i:a_i=v}\bar w_i.
\]

\noindent\textbf{Proposition 1 (set geometry and exact coverage event).}\quad
For every query \(x\) and nonnegative thresholds \(q_1\leq q_2\),
\[
  C_{q_1}(x)\subseteq C_{q_2}(x),
  \qquad C_{q_1}(x)\neq\varnothing.
\]
Moreover, for an observed label \(Y_t\),
\[
  Y_t\in C_{q_t}(X_t)
  \quad\Longleftrightarrow\quad A_t\leq q_t,
\]
and hence the conditional label-marginal coverage is exactly \(F_t(q_t)\).

\emph{Proof.}
The first claim follows because the sublevel set of \(a(x,y)\) can only grow as
its threshold increases.  The top-scoring candidate has margin zero, so it is
included whenever the threshold is nonnegative.  The final equivalence is the
definition of \(C_q(x)\); taking conditional probability gives \(F_t(q_t)\).
\hfill\(\square\)

This proposition also fixes the interpretation of the primary metric.  It is
coverage of a randomly indexed recorded answer, not simultaneous coverage of
all answers in \(Y_t(x)\).  Full-set coverage and partial answer recall are
therefore reported separately for multi-answer queries in the empirical
diagnostics.

\noindent\textbf{Proposition 2 (deterministic coverage-deviation bound).}\quad
For any finite threshold selected by the inclusive empirical rule above,
\[
  1-\alpha-\delta_t-\epsilon_t
  \;\leq\;F_t(q_t)\;\leq\;
  1-\alpha+\delta_t+\epsilon_t+\eta_t,
  \tag{1}\label{eq:coverage-decomposition}
\]
with either side clipped to \([0,1]\) when necessary.

\emph{Proof.}
By construction, \(\widehat H_{t,w}(q_t)\geq p\).  Therefore
\(H_{t,w}(q_t)\geq p-\epsilon_t\), and then
\(F_t(q_t)\geq p-\epsilon_t-\delta_t\).  Minimality of the inclusive quantile
gives \(\widehat H_{t,w}(q_t^-)<p\).  The empirical jump at \(q_t\) is the
total normalized weight of all calibration scores tied there and is at most
\(\eta_t\).  Thus
\(\widehat H_{t,w}(q_t)\leq p+\eta_t\), which implies
\(H_{t,w}(q_t)\leq p+\eta_t+\epsilon_t\) and finally the stated upper bound.
\hfill\(\square\)

Equation~\eqref{eq:coverage-decomposition} is a decomposition, not a
distribution-free guarantee.  Recency weighting can reduce \(\delta_t\) when
the recent past better matches the present, but concentrating the weights can
increase \(\epsilon_t\) and \(\eta_t\).  Conversely, long histories may reduce
sampling variability while increasing temporal mismatch.  The proposed drift
features are therefore empirical proxies for choosing this bias--variance
trade-off; they do not certify a bound on \(\delta_t\).

\subsection{Effective Sample Size and Half-Life Limits}

A transparent diagnostic for weight concentration is Kish's effective sample
size
\[
  N_{\mathrm{eff}}
  =\frac{(\sum_i w_i)^2}{\sum_i w_i^2}
  =\frac{1}{\sum_i\bar w_i^2}.
\]
For one score at each of \(W\) equally spaced time units and geometric weights
\(w_j=r^j\), \(j=0,\ldots,W-1\), where \(r=2^{-1/h}\), direct summation gives
\[
  N_{\mathrm{eff}}(h,W)
  =\frac{1+r}{1-r}\frac{1-r^W}{1+r^W}.
  \tag{2}\label{eq:neff-geometric}
\]
For a long regular stream, the limiting values are approximately 20.21, 40.40,
and 86.57 for half-lives 7, 14, and 30, respectively.  These numbers are
illustrative rather than the experiment's actual effective sample sizes: the
implementation weights every query--label score, multiple scores share a
timestamp, batch sizes vary, and the rolling pool truncates by score count.
The artifact therefore records \(N_{\mathrm{eff}}\) from the realized weights
whenever it is used for empirical interpretation.

As \(h\rightarrow\infty\), \(r\rightarrow1\) and the weighted rule approaches
equal weighting within the same rolling pool.  It need not exactly equal the
named rolling conformal baseline because that baseline uses the
finite-sample-corrected order statistic while the implemented weighted rule uses the
uncorrected level \(1-\alpha\).  As \(h\rightarrow0\), weight concentrates on
the newest timestamp represented in the pool; all scores sharing that newest
timestamp retain equal age and can retain nonzero mass.  Thus ``short
half-life'' does not mean a single calibration observation in a batched TKG.

\noindent\textbf{Corollary (ranking invariance).}\quad
For a fixed scorer and test queries, changing \(q_t\) changes set membership but
not the ordering of candidates.  MRR and Hits@\(K\) are consequently identical
for all post-hoc calibration rules applied to the same scored condition.
Calibration can trade coverage against set size or abstention, but it cannot
repair ranking damage caused by deleted training history.  Ranking metrics are
therefore reported once per scorer--deletion--seed condition rather than as if
they were improvements attributable to a post-hoc calibration rule.

Finally, formal weighted conformal results require additional structure, such
as weighted exchangeability or valid likelihood-ratio weights, and formal
non-exchangeable constructions may assign mass to the test point at
\(+\infty\).  Equation~\eqref{eq:coverage-decomposition} does not replace
those results.  It states the precise deterministic error terms for the
implemented empirical quantile and makes clear which terms the diagnostic
experiments must measure or stress.
